\section{Construcción de redes neuronales artificiales}

\begin{frame}
    El elemento básico es la neurona artificial y el nombre se basa en
    el funcionamiento de las neuronas en el cerebro:

    \pause

    \begin{definition}[Neurona artificial]
        Sea $w\in\mathbb{R}^{n}$ y $b\in\mathbb{R}$, así como
        $f\colon\mathbb{R}^{n}\to\mathbb{R}$.
        Entonces la función se llama
        \vspace*{-.5\baselineskip}\setlength\belowdisplayshortskip{0pt}
        \begin{equation*}
            f\colon\mathbb{R}^{n}\to\mathbb{R},\quad
            f\left(x\right)=\sigma\left(\sum_{i=1}^{n}w_{i}x_{i}+b\right),
        \end{equation*}
        se denomina \alert{neurona artificial}.
        El vector $w\in\mathbb{R}^{n}$ se denomina peso y
        $b\in\mathbb{R}^{n}$ es el sesgo.
        La función $\sigma\colon\mathbb{R}\to\mathbb{R}$ se denomina
        \alert{función de activación}.
    \end{definition}

    \pause

    No existe una definición clara del término función de activación.
    En cambio, se refiere a una función (no lineal) que determina
    cuándo se activa la salida de una neurona artificial.

    La variante más simple es la función de Heaviside:
    \vspace*{-.5\baselineskip}\setlength\belowdisplayshortskip{0pt}
    \begin{equation*}
        \sigma_{\text{HS}}\left(s\right)=
        H\left(s\right)\coloneqq
        \begin{cases}
            0, & s< 0.   \\
            1, & s\geq0.
        \end{cases}
    \end{equation*}

    Una función de activación continua es la
    \alert{Unidad Lineal Rectificada} (ReLU)
    \begin{equation*}
        \sigma_{\text{ReLU}}\left(s\right)=
        \operatorname{ReLu}\left(s\right)\coloneqq
        \begin{cases}
            0, & s\leq 0. \\
            s, & s>0.
        \end{cases}
    \end{equation*}

    Las \alert{activaciones sigmoideas} son continuamente
    diferenciables (con una frecuencia arbitraria).
    \begin{equation*}
        \sigma\left(s\right)=
        \frac{1}{1+\exp\left(-s\right)},
    \end{equation*}
    o la tangente hiperbólica
    \begin{equation*}
        \sigma_{\tanh}\left(s\right)=
        \tanh\left(s\right).
    \end{equation*}

    Estas dos están estrechamente relacionadas, es
    \vspace*{-.5\baselineskip}\setlength\belowdisplayshortskip{0pt}
    \begin{equation*}
        \sigma\left(s\right)=
        \frac{1}{2}
        \tanh\left(\frac{x}{2}\right)+
        \frac{1}{2}.
    \end{equation*}
\end{frame}

\begin{frame}
    \begin{figure}[ht!]
        \centering
        \includegraphics[width=.7\paperwidth]{1110}
        \caption{Izquierda: Esquema de una neurona artificial con tres
            señales de entrada, la capa lineal y la activación no lineal.
            Derecha: El desarrollo de algunas funciones de activación
            típicas.}
        \label{fig:11.10}
    \end{figure}
    En la Figura~\ref{fig:11.10} mostramos el curso de estas funciones
    de activación así como la estructura esquemática de una neurona
    artificial.

    \pause

    \begin{example}[Clasificación con neuronas artificiales]
        Una neurona artificial con activación de Heaviside
        \vspace*{-.5\baselineskip}\setlength\belowdisplayshortskip{0pt}
        \begin{equation*}
            f\left(x\right)=
            H\left({\left\langle w,x\right\rangle}_{2}+b\right),
        \end{equation*}
        divide el espacio en un hiperplano, porque cada plano con vector
        normal $n$ y punto seleccionado $x_{0}$ está dado por
        \begin{equation*}
            0=n^{T}\left(x-x_{0}\right)=
            n^{T}x-
            n^{T}x_{0}.
        \end{equation*}

        \pause

        Por lo tanto, el peso corresponde al vector normal, y el sesgo es
        el producto escalar con el punto seleccionado.
        Una neurona divide así el espacio en dos mitades más allá del
        plano, y esta es una de las tareas fundamentales para las que se
        pueden utilizar redes neuronales artificiales.
    \end{example}
\end{frame}

\begin{frame}
    Dados $x_{1},\dotsc,x_{N}$, a cada una se le asigna una propiedad
    $y_{i}$, simplemente $y_{i}\in\left\{0,1\right\}$.
    Buscamos los pesos $w$ y el sesgo $b$ de la neurona artificial
    tales que $f\left[w,b\right](x_{i})=y_{i}$.
    Es decir, buscamos un modelo que realice la clasificación de los
    datos en función de las propiedades.
    Los pesos refuerzan o debilitan la influencia de la entrada y, por
    lo tanto, la transición de la no activación de una neurona a la
    activación.
    El sesgo asegura un cambio en la activación.
    Un modelo compuesto por una sola neurona es muy limitado, y debe
    asumirse que la tarea no puede resolverse con exactitud.
    En su lugar, intentamos encontrar el mejor modelo; es decir,
    consideramos el problema de minimización de forma análoga a la
    regresión de la Sección 4.5 y el Excurso 11.4.
    \begin{equation*}
        \min_{\left(w,b\right)}
        \sum_{i=1}^{N}
        {\left\|y_{i}-f\left[w,b\right]\left(x_{i}\right)\right\|}^{2}.
    \end{equation*}

    \pause

    Las redes neuronales artificiales (RNA) se componen de neuronas
    artificiales individuales.
    Las RNA más simples son las llamadas redes de una sola capa:

    \begin{definition}[Red de una sola capa]
        Sea $W\in\mathbb{R}^{n\times m}$ y $b\in\mathbb{R}^{n}$, así como
        $\sigma\colon\mathbb{R}\to\mathbb{R}$ una función de activación.
        \vspace*{-.5\baselineskip}\setlength\belowdisplayshortskip{0pt}
        \begin{equation*}
            f\left(x\right)=
            \sigma\left(Wx+b\right),
        \end{equation*}
        En ese caso, la función se denomina red neuronal artificial con
        una capa y $n$ neuronas artificiales.
        En este caso, la aplicación de la función de activación debe
        entenderse por componentes, es decir, para $y\in\mathbb{R}^{n}$
        es
        \begin{equation*}
            \sigma\left(y\right)=
            {\begin{pmatrix}
                \sigma\left(y_{i}\right)
            \end{pmatrix}}_{i=1,\dotsc,n}.
        \end{equation*}
    \end{definition}
    A partir de esta definición, pasamos rápidamente a redes neuronales
    artificiales más generales, las llamadas redes neuronales
    profundas.
    Para ello, conectamos varias redes de una sola capa en serie.
    La salida de la capa $l$ es la entrada de la capa $l+1$.
\end{frame}

\begin{frame}
    \begin{definition}[Red Neuronal Profunda]
        Sea $L\in\mathbb{N}$ el número de capas.
        Para $l=1,\dotsc,L$, sea
        $W^{l}\in\mathbb{N}^{n_{l}\times n_{l-1}}$ y
        $b^{l}\in\mathbb{R}^{n_{l}}$ el número de capas.
        La función de activación viene dada por $\sigma\colon\mathbb{R}\to\mathbb{R}$.
        Entonces, $f\colon\mathbb{R}^{n_{0}}\to\mathbb{R}^{n_{L}}$ viene dada por
        \vspace*{-.5\baselineskip}\setlength\belowdisplayshortskip{0pt}
        \begin{align*}
            x^{0}           & =x.                                                  \\
            \forall l=1,\dotsc,L-1:
            z^{l}           & =W^{l}x^{l-1}+b^{l}.                                 \\
            \forall l=1,\dotsc,L-1:
            x^{l}           & =f^{l}\left(x^{l-1}\right)=\sigma\left(z^{l}\right). \\
            f\left(x\right) & =x^{L}=W^{L}x^{L-1}+b^{L}.
        \end{align*}
    \end{definition}
    una \alert{red neuronal profunda}.
    La primera capa $f^{1}\left(\cdot\right)$ se denomina capa de
    entrada, la última capa $f^{L}\left(\cdot\right)$ es la capa de
    salida y, por lo general, no se añade ninguna función de activación.
    Las capas internas se denominan capas ocultas.
    Nos referimos a la arquitectura de una red neuronal profunda como
    $\mathcal{N}\left(\sigma;n_{0},n_{1},\dotsc,n_{L}\right)$ o, para
    abreviar, como $\mathcal{L}\left(L,n\right)$ con $n=\max_{l}n_{l}$.
    Denotamos el espacio de parámetros correspondiente por
    \begin{equation*}
        \mathcal{P}_{\mathcal{N}}=
        \left\{
        W=\left(W_{1},\dotsc,W_{L}\right),b=\left(b_{1},\dotsc,b_{L}\right),
        W_{l}\in\mathbb{R}^{n_{l}\times n_{l-1}},b_{l}\in\mathbb{R}^{n_{l}}
        \right\}.
    \end{equation*}
    Para una elección de parámetro $\left(W,b\right)\in\mathcal{P}$,
    denotamos por $f\colon\left[W,b\right]\colon\mathbb{R}^{n_{0}}\to\mathbb{R}^{n_{L}}$
    la función resultante según (11.18).
    El conjunto de realizaciones, es decir, el conjunto de funciones
    que se pueden representar en la arquitectura, se denota por
    \begin{equation*}
        \mathcal{F}_{\mathcal{N}}=
        \left\{f\left[W,b\right]\mid \left(W,b\right)\in\mathcal{P}_{\mathcal{N}}\right\}.
    \end{equation*}
    Contando los parámetros libres se obtiene:
    \begin{corollary}[Número de parámetros libres]
        Sea $\mathcal{N}\left(\sigma;n_{0},\dotsc,n_{L}\right)$ una red
        neuronal artificial con $L$ capas y $n_{l}$ neuronas artificiales
        cada una.
        Entonces, el espacio de parámetros $\mathcal{P}_{\mathcal{N}}$
        \vspace*{-.5\baselineskip}\setlength\belowdisplayshortskip{0pt}
        \begin{equation*}
            \#\mathcal{P}_{\mathcal{N}}=
            \sum_{l=1}^{L}
            n_{l}\left(n_{l-1}+1\right)
        \end{equation*}
        tiene un total de parámetros libres, de los cuales
        $\sum_{l=1}^{L}n_{l}n_{l-1}$ pesos y $\sum_{l=1}^{L}n_{l}$
        valores de sesgo.
    \end{corollary}
\end{frame}

\begin{frame}
    Una operación básica en aprendizaje automático es la evaluación de
    una red dada $\mathcal{N}\left(\sigma;n_{0},\dotsc,n_{L}\right)$
    para parámetros dados
    $\left(W,b\right)\in\mathcal{P}_{\mathcal{N}}$ para un punto de
    datos $x\in\mathbb{R}^{n_{0}}$ o, más frecuentemente, para una
    lista completa de (posiblemente muchos) puntos de datos
    $x_{i}\in\mathbb{R}^{n_{0}}$ para $i=1,\dotsc,N_{tr}$.

    \begin{algorithm}[H]
        \TitleOfAlgo{Evaluación de una red neuronal artificial}
        \Entrada{Red neuronal artificial
        $\mathcal{N}\left(\sigma;n_{0},\dotsc,n_{L}\right)$,
        Parámetro $\left(W,b\right)\in\mathcal{P}$ y entrada
        $x\in\mathbb{R}^{n_{0}}$.}
        $x^{0}\leftarrow x$\;
        \ParaCada{$l=1,\dotsc,L$}{$x^{l}\leftarrow W^{l}x^{l-1}+b^{l}$\;}
        \eSSi{\normalfont $l<L$}{$x^{l}\leftarrow\sigma\left(z^{l}\right)$\;}{$x^{l}\leftarrow z^{l}$\;}
        \Resultado{Salida de red $x^{L}=f\left[W,b\right]\left(x\right)$.}
        % \caption{}\label{alg:evalknn}
    \end{algorithm}

    \begin{alertblock}{Arquitecturas de Red}
        La definición anterior introduce solo el tipo más simple de redes
        neuronales artificiales: todas las salidas de la capa $l-1$ son
        entradas a todas las neuronas artificiales de la capa $l$.
        Estas redes también se denominan \alert{perceptrones multicapa}.
        Pertenecen a la categoría de
        \alert{redes neuronales de propagación hacia adelante}, es
        decir, la información fluye en una sola
        dirección, de la capa $1$ a la capa $2$, y así sucesivamente.
        Las redes neuronales en las que la salida de la capa $l$ también
        puede servir como entrada de capas anteriores $l^{\prime}<l$ se
        denominan \alert{redes neuronales recurrentes}.
    \end{alertblock}

    \pause

    En principio, se puede utilizar una función de activación
    diferente en cada capa.
    Si se aplica una función de activación en la última capa, se
    restringe el espacio de la imagen a $\left(0,1\right)$ en el caso
    de la activación sigmoidea y a los números no negativos en el
    caso de la activación $\operatorname{ReLU}$.
\end{frame}

\begin{frame}
    Las llamadas conexiones de salto omiten capas individuales.
    Esto significa que la salida de la capa $l$ puede utilizarse
    directamente como entrada, por ejemplo, para la capa $l+2$.
    Otra modificación sencilla son las \alert{redes residuales}, en
    las que se forma una capa con la forma
    \begin{equation*}
        f\left(x\right)=
        x+
        \sigma\left(Wx+b\right),
    \end{equation*}
    es decir, la salida de la neurona artificial se añade a la
    entrada.

    Las \alert{redes neuronales convolucionales} desempeñan un papel
    fundamental en el procesamiento de imágenes.
    En este caso, cada \alert{capa} es una convolución.
    La particularidad de las redes convolucionales es que el tamaño de
    la entrada $x$ no tiene por qué estar necesariamente vinculado al
    número de pesos.
    Con los pesos de convolución $W\in\mathbb{R}^{f}$ definimos
    \begin{equation*}
        f_{\text{conv}}\left(x\right)=
        \sigma\left(W\ast x\right),\qquad
        {\left(W\ast x\right)}_{i}=
        \sum_{j=1}^{f}
        W_{j}x_{i-j},
    \end{equation*}
    donde $x_{0}=x_{-1}=\cdots=x_{1-f}1=0$.
    Los mismos pesos se aplican siempre solo a una pequeña parte de
    las entradas.
    De esta manera, las funciones en espacios de dimensiones
    arbitrarias pueden representarse con solo unos pocos parámetros
    libres (los pesos).

    Las arquitecturas modernas, como las utilizadas en modelos de
    lenguaje como ChatGPT, son mucho más complejas, pero también
    utilizan y combinan los modelos básicos presentados aquí.

    Las redes neuronales artificiales utilizadas en la práctica son
    enormemente grandes; por ejemplo, ChatGPT-3 tiene aproximadamente
    $200$ mil millones, o $2\cdot 10^{11}$ de parámetros libres.
\end{frame}

\begin{frame}
    En la Figura~\ref{fig:11.11} representamos gráficamente algunas
    redes neuronales artificiales.
    Se puede encontrar una descripción detallada en la literatura
        [11, 17, 43, 51, 66].
    \begin{figure}[ht!]
        \centering
        \includegraphics[width=.7\paperwidth]{1111}
        \caption{Ejemplos de redes neuronales artificiales.
            Arriba a la izquierda: perceptrón multicapa simple
            $\mathcal{N}\left(\sigma;2,3,3,2\right)$ con activación
            sigmoidea (excepto la capa de salida).
            Arriba a la derecha: esta red representa las funciones
            $f\colon\mathbb{R}^{2}\to\mathbb{R}^{2}$.
            Red neuronal profunda con dos entradas y salidas y $L-1$
            capas ocultas, cada una con $n$ neuronas.
            Abajo a la izquierda: red neuronal recurrente, donde las
            salidas también pueden ser entradasde capas anteriores.
            Abajo a la derecha: red neuronal residual artificial.
            La entrada se suma a la salida.
            Esta red representa las funciones
            $f\colon\mathbb{R}\to\mathbb{R}$
            con la estructura
            $f\left(x\right)=x+f\left[W,b\right]\left(x\right)$.}
        \label{fig:11.11}
    \end{figure}
\end{frame}

\begin{frame}
    Consideramos únicamente redes neuronales artificiales profundas
    del tipo definido en la Definición 11.14 y definimos una red con
    $L$ capas como una función $f\colon\mathbb{R}^{n_{0}}\to\mathbb{R}^{n_{L}}$
    con las matrices de pesos $W_{l}$ y los vectores de sesgo $b_{l}$ como
    parámetros libres.
    Primero, establecemos una conexión simple pero importante:
    \begin{theorem}
        Sea
        \begin{math}
            \mathcal{N}=
            \mathcal{N}\left(\sigma,n_{0},\dotsc,n_{L}\right)
        \end{math}
        una arquitectura de red con $L$ capas y parámetros
        $W_{l}\in\mathbb{R}^{n_{l}\times n_{l-1}}$
        y $b^{l}\in\mathbb{R}^{n_{l}}$ para $l=1,\dotsc,L$.
        El conjunto de realizaciones
        \begin{equation*}
            \mathcal{F}_{\mathcal{N}}\coloneqq
            \left\{
            f\left[W,b\right]\colon\mathbb{R}^{n_{0}}\to\mathbb{R}^{n_{L}}\mid
            \left(W,b\right)\in\mathcal{P}_{\mathcal{N}}
            \right\}
        \end{equation*}
        no es generalmente un espacio vectorial.
        En particular, $\mathcal{F}_{\mathcal{N}}$ no es cerrado con respecto
        a la suma ni a la multiplicación escalar.
    \end{theorem}
\end{frame}